\documentclass[a4paper,12pt]{article}
\usepackage[utf8]{inputenc}
\usepackage[T1]{fontenc}
\usepackage[ngerman]{babel}
\usepackage{amsmath, amssymb}
\usepackage{geometry}
\usepackage{tikz}
\usetikzlibrary{shapes.geometric, arrows.meta, positioning, fit, backgrounds}
\usepackage{parskip}
\usepackage{titlesec}
\usepackage{booktabs}
\usepackage{enumitem}

\geometry{left=3cm, right=3cm, top=2.5cm, bottom=2.5cm}

\titleformat{\section}{\large\bfseries}{}{0em}{}[\titlerule]
\titleformat{\subsection}{\normalsize\bfseries}{}{0em}{}

\tikzset{
  box/.style={rectangle, draw, rounded corners=3pt, text width=4.2cm,
              align=center, minimum height=0.9cm, font=\small, inner sep=4pt},
  boxw/.style={box, text width=5.2cm},
  boxn/.style={box, text width=3.8cm},
  arrow/.style={-{Latex[length=2mm]}, thick},
}

\begin{document}

\begin{center}
  {\LARGE\bfseries Studierhinweise zu Lektion 3}\\[0.5em]
  {\large Analysis (Modul 61211)}
\end{center}

\vspace{1em}

In dieser Kurseinheit steht der Begriff der Differenzierbarkeit im Mittelpunkt.
Zur Vorbereitung seiner Definition (in Abschnitt~3.5) werden zwei Themen
behandelt: der Begriff des Grenzwerts einer Funktion (er beruht auf dem
Stetigkeitsbegriff) und der Begriff der Matrix (Matrizen dienen der Darstellung
linearer Funktionen in $\mathrm{Abb}(\mathbb{R}^n, \mathbb{R}^m)$).

Abschnitt~3.1 geht zuru\"uck auf das, was bereits in MG \"uber Grenzwerte von
Funktionen erfahren wurde; hinzugekommen sind Grenzwerte in den uneigentlichen
Punkten $-\infty$ und $\infty$ sowie uneigentliche Grenzwerte von Funktionen.
Im zweiten Abschnitt wird der Grenzwert f\"ur allgemeinere Funktionen definiert.
Die Abschnitte~3.3 und~3.4 sind erneut R\"uckblicke auf schon Bekanntes aus MG:
Differenzierbarkeit von Funktionen (unter dem Aspekt der guten Approximierbarkeit
durch affine Funktionen) und Fakten \"uber lineare Funktionen und Matrizen.
Den zentralen Abschnitt~3.5 \"uber den Differenzierbarkeitsbegriff und seine
elementaren Eigenschaften sollten Sie besonders sorgf\"altig durcharbeiten.
Abschnitt~3.6 stellt die Hilfsmittel bereit, Ableitungen bequem auszurechnen.

% -----------------------------------------------------------------------
\section*{Struktur der Kurseinheit 3}

\begin{center}
\begin{tikzpicture}[node distance=0.5cm and 0.6cm]

  % Column headers
  \node[font=\bfseries\small] (hA) at (0,0)    {R\"uckblick auf MG};
  \node[font=\bfseries\small] (hB) at (6.0,0)  {Erg\"anzungen zu MG};
  \node[font=\bfseries\small] (hC) at (12.0,0) {$f \in \mathrm{Abb}(M,Y)$, $M \subseteq X$};

  % ---- Left column: MG-Rückblick ----
  \node[box, below=0.8cm of hA] (l1) {3.1 stetige Fortsetzung};
  \node[box, below=0.4cm of l1] (l2) {3.1 Grenzwerte\\von Funktionen};
  \node[box, below=0.4cm of l2] (l3) {3.1 Operieren\\mit Grenzwerten};
  \node[box, below=0.8cm of l3] (l4) {3.3 Differenzierbarkeit,\\Ableitung};
  \node[box, below=0.4cm of l4] (l5) {3.3 Differenzierbarkeit\\und Stetigkeit};
  \node[box, below=0.4cm of l5] (l6) {3.3 Differenziations-\\regeln};
  \node[box, below=0.4cm of l6] (l7) {3.3 1. und 2.\\Mittelwertsatz};

  % ---- Middle column: Ergänzungen ----
  \node[box, below=0.8cm of hB] (m1) {3.1 Grenzwerte in\\$\infty$ oder in $-\infty$,\\uneigentliche Grenzwerte};
  \node[box, below=0.4cm of m1] (m2) {3.4 lin.\ Abbildungen\\in $\mathrm{Abb}(\mathbb{R}^n, \mathbb{R}^m)$\\und ihre Matrix};
  \node[box, below=0.4cm of m2] (m3) {3.4 $\mathrm{Hom}(\mathbb{R}^n,\mathbb{R}^m)$\\als normierter Raum};
  \node[box, below=0.4cm of m3] (m4) {3.3 Regeln von\\de l'Hospital};

  % ---- Right column: new content ----
  \node[box, below=0.8cm of hC] (r1) {3.2 stetige Fortsetzung};
  \node[box, below=0.4cm of r1] (r2) {3.2 Grenzwerte\\von Funktionen};
  \node[box, below=0.4cm of r2] (r3) {3.2 Operieren\\mit Grenzwerten};
  \node[box, below=0.4cm of r3] (r4) {3.5 (totale)\\Differenzierbarkeit,\\Ableitung};
  \node[box, below=0.4cm of r4] (r5) {3.5 Differenzierbarkeit\\und Stetigkeit};
  \node[box, below=0.4cm of r5] (r6) {3.5 Differenziations-\\regeln};
  \node[box, below=0.4cm of r6] (r7) {3.5 Mittelwertsatz};
  \node[box, below=0.4cm of r7] (r8) {3.6 partielle\\Differenzierbarkeit,\\partielle Ableitungen};
  \node[box, below=0.4cm of r8] (r9) {3.6 Richtungs-\\ableitung};
  \node[box, below=0.4cm of r9] (r10){3.6 partielle und totale\\Differenzierbarkeit,\\Funktionalmatrix};

  % Arrows left column
  \draw[arrow] (l1) -- (l2);
  \draw[arrow] (l2) -- (l3);
  \draw[arrow] (l3) -- (l4);
  \draw[arrow] (l4) -- (l5);
  \draw[arrow] (l5) -- (l6);
  \draw[arrow] (l6) -- (l7);

  % Arrows right column
  \draw[arrow] (r1) -- (r2);
  \draw[arrow] (r2) -- (r3);
  \draw[arrow] (r3) -- (r4);
  \draw[arrow] (r4) -- (r5);
  \draw[arrow] (r5) -- (r6);
  \draw[arrow] (r6) -- (r7);
  \draw[arrow] (r7) -- (r8);
  \draw[arrow] (r8) -- (r9);
  \draw[arrow] (r9) -- (r10);

  % Cross arrows
  \draw[arrow, dashed] (l3.east) -- ++(0.2,0) |- (r1.west);
  \draw[arrow, dashed] (m2.east) -- ++(0.2,0) |- (r4.west);
  \draw[arrow, dashed] (m3.east) -- ++(0.2,0) |- (r4.west);

\end{tikzpicture}
\end{center}

\newpage

% -----------------------------------------------------------------------
\section*{Zielelement 3.1 -- R\"uckblick und Erg\"anzungen: Grenzwerte reeller Funktionen}

\subsection*{Lerninhalte}

\begin{center}
\begin{tikzpicture}[node distance=0.5cm and 1.0cm]
  \node[box] (sf)   {stetige Fortsetzung};
  \node[box, below=0.5cm of sf]   (gw)  {Grenzwert\\einer Funktion};
  \node[box, below=0.5cm of gw]   (kri) {Umgebungskriterium,\\$\varepsilon$-$\delta$-Kriterium,\\Folgenkriterium};
  \node[box, right=1.0cm of kri]  (op)  {Regeln f\"ur das\\Operieren mit Grenzwerten};
  \node[box, below=0.5cm of kri]  (cau) {Cauchykriterium};
  \node[box, below=0.5cm of cau]  (inf) {Grenzwert in $-\infty$\\und in $\infty$};
  \node[box, below=0.5cm of inf]  (uei) {uneigentlicher\\Grenzwert};
  \draw[arrow] (sf) -- (gw);
  \draw[arrow] (gw) -- (kri);
  \draw[arrow] (op.south) |- (cau.east);
  \draw[arrow] (kri) -- (cau);
  \draw[arrow] (cau) -- (inf);
  \draw[arrow] (inf) -- (uei);
\end{tikzpicture}
\end{center}

\subsection*{Lernziele}

Nach Durcharbeiten dieses Abschnitts sollten Sie:
\begin{itemize}
  \item sich die Begriffe ,,stetige Fortsetzung'' und ,,Grenzwert einer Funktion''
        wieder vergegenw\"artigt haben und miteinander in Verbindung bringen
        k\"onnen,
  \item die Regeln f\"ur das Operieren mit Grenzwerten angeben und anwenden
        k\"onnen,
  \item die Zeichen $\lim_{x \to \infty} f(x) = a$ und $\lim_{x \to a} f(x) = \infty$
        erl\"autern k\"onnen.
\end{itemize}

\subsection*{Selbstkontrollelement 3.1}

Sei $(a_k)$ eine konvergente Folge mit Grenzwert $a$, also
$\lim_{k\to\infty} a_k = a$. ,,Sehen'' Sie, dass die Funktion
$f : \bigl\{\tfrac{1}{k} \mid k \in \mathbb{N}\bigr\} \to \mathbb{R}$,
$\tfrac{1}{k} \mapsto f(\tfrac{1}{k}) := a_k$ einen Grenzwert in $0$ hat?
Welchen? Geben Sie eine pr\"azise Begr\"undung.

\newpage

% -----------------------------------------------------------------------
\section*{Zielelement 3.2 -- Grenzwerte von Funktionen auf normierten R\"aumen}

\subsection*{Lerninhalte}

\begin{center}
\begin{tikzpicture}[node distance=0.5cm and 1.0cm]
  \node[box] (sf)  {stetige Fortsetzung,\\Grenzwert einer Funktion};
  \node[box, right=1.0cm of sf] (kri) {Umgebungskriterium,\\$\varepsilon$-$\delta$-Kriterium,\\Folgenkriterium};
  \node[box, below=0.5cm of sf]  (res) {Restriktion und\\Verkettung};
  \node[box, below=0.5cm of res] (op)  {Regeln f\"ur das\\Operieren mit Grenzwerten};
  \draw[arrow] (sf) -- (res);
  \draw[arrow] (kri.south) |- (res.east);
  \draw[arrow] (res) -- (op);
\end{tikzpicture}
\end{center}

\subsection*{Lernziele}

Nach Durcharbeiten dieses Abschnitts sollten Sie in der Lage sein,
\begin{itemize}
  \item die Begriffe ,,stetige Fortsetzung'' und ,,Grenzwert einer Funktion''
        auf normierten R\"aumen zu definieren und mit dem Begriff der
        Stetigkeit in Verbindung zu bringen,
  \item Regeln f\"ur das Operieren mit Grenzwerten anzugeben und anzuwenden.
\end{itemize}

\subsection*{Selbstkontrollelement 3.2}

Sei $a \in M \subseteq \mathbb{R}^n$, $a$ ein H\"aufungspunkt von $M$, und
f\"ur $f : M \to \mathbb{R}^m$ existiere $b := \lim_{x \to a} f(x)$.
\begin{itemize}
  \item Kann $b \neq f(a)$ sein?
  \item Existiert f\"ur $\tilde{f} := f|_{M \setminus \{a\}}$ der Grenzwert
        $\lim_{x \to a} \tilde{f}(x)$?
  \item \"Andern sich die Antworten, wenn $f$ als stetig in $a$ vorausgesetzt wird?
\end{itemize}

\newpage

% -----------------------------------------------------------------------
\section*{Zielelement 3.3 -- R\"uckblick und Erg\"anzungen: Differenzierbarkeit in $\mathbb{R}$}

\subsection*{Lerninhalte}

\begin{center}
\begin{tikzpicture}[node distance=0.5cm and 1.0cm]
  \node[box] (diff)  {Differenzierbarkeit,\\Begriff der Ableitung};
  \node[box, right=1.0cm of diff] (geo) {geometrische und\\physikalische Interpretation};
  \node[box, below=0.5cm of diff] (ds)   {Differenzierbarkeit\\und Stetigkeit};
  \node[box, below=0.5cm of ds]   (regeln) {Differenziationsregeln};
  \node[box, right=1.0cm of regeln] (umk) {Differenziation der\\Umkehrfunktion};
  \node[box, below=0.5cm of regeln] (mws) {Mittelwerts\"atze};
  \node[box, below=0.5cm of mws]  (hos)  {Regeln von\\de l'Hospital};
  \node[box, below=0.5cm of hos]  (bsp)  {Beispiele: Polynomfunktion,\\rationale Funktion, $p$-te Wurzel};
  \draw[arrow] (diff) -- (ds);
  \draw[arrow] (geo.south) |- (ds.east);
  \draw[arrow] (ds) -- (regeln);
  \draw[arrow] (umk.south) |- (mws.east);
  \draw[arrow] (regeln) -- (mws);
  \draw[arrow] (mws) -- (hos);
  \draw[arrow] (hos) -- (bsp);
\end{tikzpicture}
\end{center}

\subsection*{Lernziele}

Nach Durcharbeiten dieses Abschnitts sollten Sie:
\begin{itemize}
  \item sich die Begriffe ,,Differenzierbarkeit'' und ,,Ableitung'' wieder
        vergegenw\"artigt haben,
  \item die Beziehung zwischen Stetigkeit und Differenzierbarkeit kennen,
  \item die wichtigsten Differenziationsregeln angeben und anwenden k\"onnen,
  \item wenigstens eine der Regeln von de l'Hospital (auch f\"ur Grenzwerte in
        $\infty$ oder $-\infty$ sowie f\"ur uneigentliche Grenzwerte)
        erl\"autern k\"onnen.
\end{itemize}

\subsection*{Selbstkontrollelement 3.3}

Sei $f : \mathbb{R} \to \mathbb{R}$ eine Funktion mit $|f(x)| \leq x^2$ f\"ur
alle $x \in \mathbb{R}$. ,,Sehen'' Sie, dass ein solches $f$ in $a := 0$
differenzierbar ist? Mit welcher Ableitung? (Betrachten Sie
$\tfrac{f(x) - f(0)}{x - 0}$.)

\newpage

% -----------------------------------------------------------------------
\section*{Zielelement 3.4 -- R\"uckblick und Erg\"anzungen: Der Raum $\mathrm{Hom}(\mathbb{R}^n, \mathbb{R}^m)$}

\subsection*{Lerninhalte}

\begin{center}
\begin{tikzpicture}[node distance=0.5cm and 1.0cm]
  \node[box] (lin)  {lineare Abbildungen\\in $\mathrm{Abb}(\mathbb{R}^n, \mathbb{R}^m)$\\und ihre Matrix};
  \node[box, right=1.0cm of lin] (op) {Operationen mit\\linearen Abbildungen,\\Matrizen};
  \node[box, below=0.5cm of lin] (hom) {$\mathrm{Hom}(\mathbb{R}^n, \mathbb{R}^m)$\\als normierter Raum};
  \draw[arrow] (lin) -- (hom);
  \draw[arrow] (op.south) |- (hom.east);
\end{tikzpicture}
\end{center}

\subsection*{Lernziele}

Nach Durcharbeiten dieses Abschnitts sollten Sie in der Lage sein,
\begin{itemize}
  \item lineare Abbildungen von $\mathbb{R}^n$ nach $\mathbb{R}^m$ mithilfe
        ihrer Matrix auszudr\"ucken (die kanonischen Basen von $\mathbb{R}^n$
        und $\mathbb{R}^m$ vorausgesetzt),
  \item das Operieren mit linearen Abbildungen dem mit ihren Matrizen
        gegen\"uberzustellen, insbesondere die Verkettung linearer Abbildungen
        und das Matrizenprodukt zu bilden und miteinander in Beziehung zu
        setzen,
  \item den Raum $\mathbb{R}^{mn}$ als reellen Vektorraum zu beschreiben und
        eine Norm auf ihm anzugeben.
\end{itemize}

\subsection*{Selbstkontrollelement 3.4}

Sei $f \in \mathrm{Hom}(\mathbb{R}^n, \mathbb{R}^m)$, $\mathrm{id}_n : \mathbb{R}^n \to \mathbb{R}^n$
die identische Abbildung auf $\mathbb{R}^n$ und $\mathrm{id}_m : \mathbb{R}^m \to \mathbb{R}^m$
die identische Abbildung auf $\mathbb{R}^m$. ,,Sehen'' Sie, wie sich die
Beziehungen $f \circ \mathrm{id}_n = f$ und $\mathrm{id}_m \circ f = f$ als
Matrixgleichungen ausdr\"ucken?

\newpage

% -----------------------------------------------------------------------
\section*{Zielelement 3.5 -- Differenzierbare Funktionen}

\subsection*{Lerninhalte}

\begin{center}
\begin{tikzpicture}[node distance=0.5cm and 1.0cm]
  \node[box] (diff)  {Differenzierbarkeit,\\Begriff der Ableitung};
  \node[box, right=1.0cm of diff] (geo) {geometrische Inter-\\pretation, Beispiele};
  \node[box, below=0.5cm of diff] (ds)  {Differenzierbarkeit\\und Stetigkeit};
  \node[box, below=0.5cm of ds]   (regeln) {Differenziationsregeln};
  \node[box, below=0.5cm of regeln] (mws) {Mittelwertsatz};
  \node[box, right=1.0cm of mws]  (szw) {Satz vom endlichen\\Zuwachs};
  \node[box, below=0.5cm of mws]  (konst) {Charakterisierung\\der konstanten Funktionen};
  \draw[arrow] (diff) -- (ds);
  \draw[arrow] (geo.south) |- (ds.east);
  \draw[arrow] (ds) -- (regeln);
  \draw[arrow] (regeln) -- (mws);
  \draw[arrow] (szw.south) |- (konst.east);
  \draw[arrow] (mws) -- (konst);
\end{tikzpicture}
\end{center}

\subsection*{Lernziele}

Nach Durcharbeiten dieses Abschnitts sollten Sie in der Lage sein,
\begin{itemize}
  \item die Begriffe ,,Differenzierbarkeit'' und ,,Ableitung'' (f\"ur
        $f \in \mathrm{Abb}(M, \mathbb{R}^m)$ mit $M \subseteq \mathbb{R}^n$)
        zu entwickeln und an Beispielen zu erl\"autern,
  \item sich die Beziehung zwischen Differenzierbarkeit und Stetigkeit klar
        zu machen,
  \item die wichtigsten Differenziationsregeln anzugeben und anzuwenden,
  \item den Mittelwertsatz zu formulieren,
  \item die konstanten Funktionen auf offenen, zusammenh\"angenden Mengen zu
        charakterisieren.
\end{itemize}

\subsection*{Selbstkontrollelement 3.5}

Sei $f : \mathbb{R}^2 \to \mathbb{R}^2$,
$x = \bigl(\begin{smallmatrix} x_1 \\ x_2 \end{smallmatrix}\bigr) \mapsto
 f(x) = f(x_1, x_2) := \bigl(\begin{smallmatrix} 3x_1 - x_2 + 2 \\ 2x_2 - 1 \end{smallmatrix}\bigr)$,
und sei $a \in \mathbb{R}^2$. Was ist $f'(a)$ und was $f'(a)x$, insbesondere
f\"ur $a := \bigl(\begin{smallmatrix} 1 \\ 0 \end{smallmatrix}\bigr)$ und
$x := \bigl(\begin{smallmatrix} 1 \\ 2 \end{smallmatrix}\bigr)$?

\newpage

% -----------------------------------------------------------------------
\section*{Zielelement 3.6 -- Partielle Ableitungen, Richtungsableitungen}

\subsection*{Lerninhalte}

\begin{center}
\begin{tikzpicture}[node distance=0.5cm and 1.0cm]
  \node[box] (pd)   {partielle Differenzierbarkeit,\\partielle Ableitung};
  \node[box, right=1.0cm of pd] (rd) {Differenzierbarkeit in\\einer Richtung,\\Richtungsableitung};
  \node[box, below=0.5cm of pd]   (tp) {Zusammenhang zwischen\\totaler und partieller\\Differenzierbarkeit};
  \node[box, right=1.0cm of tp]   (ar) {Zusammenhang zwischen\\(totaler) Ableitung und\\den Richtungsableitungen};
  \node[box, below=0.5cm of tp]   (gr) {Gradient};
  \node[box, below=0.5cm of gr]   (fm) {Funktionalmatrix};
  \draw[arrow] (pd) -- (tp);
  \draw[arrow] (rd.south) |- (ar.west);
  \draw[arrow] (rd) -- (ar);
  \draw[arrow] (tp) -- (gr);
  \draw[arrow] (ar.south) |- (gr.east);
  \draw[arrow] (gr) -- (fm);
\end{tikzpicture}
\end{center}

\subsection*{Lernziele}

Nach Durcharbeiten dieses Abschnitts sollten Sie in der Lage sein,
\begin{itemize}
  \item partielle Ableitungen und -- allgemeiner -- Richtungsableitungen zu
        definieren und in konkreten F\"allen zu berechnen,
  \item den Zusammenhang zwischen (totaler) Differenzierbarkeit und Existenz
        von partiellen Ableitungen und Richtungsableitungen herzustellen und
        die Bedeutung des Gradienten zu erl\"autern,
\end{itemize}
und Sie sollten wissen,
\begin{itemize}
  \item dass $f'(a)$ die Funktionalmatrix (aller partiellen Ableitungen) ist.
\end{itemize}

\subsection*{Selbstkontrollelement 3.6}

Welche Implikationen bestehen zwischen den folgenden Aussagen \"uber eine
Funktion $f : M \to \mathbb{R}$ mit $M \subseteq \mathbb{R}^n$ und einem
inneren Punkt $a \in M$?
\begin{enumerate}
  \item $D_k f(a)$ existiert f\"ur jedes $k \in \{1, \ldots, n\}$.
  \item $f$ ist in $a$ differenzierbar ($f'(a)$ existiert).
  \item $D_v f(a)$ existiert f\"ur jedes $v \in \mathbb{R}^n$ mit $\|v\|_2 = 1$.
  \item $f$ ist stetig in $a$.
\end{enumerate}

\end{document}
